% ~0,5 página
\section{Conclusões e trabalho futuro}
\label{sec:conclusao}

A revisão mostra que as extrapolações físicas continuam a ser referências
fortes, que os modelos clássicos têm resultados sólidos em tarefas
específicas e que um erro de previsão menor não garante, por si só, um
plano melhor. Os resultados preliminares confirmam os dois lados. Com as
variáveis certas, uma floresta aleatória reduz para metade o erro da
velocidade constante a 3 e a \un{5}{s} em cenas urbanas reais. No
planeador atual, essa melhoria ainda não se traduz em menos conflitos.

O trabalho seguinte decorre das limitações identificadas:
\begin{itemize}
  \item planear ao longo da faixa, usando o mapa, com as previsões nos
    três horizontes em vez de uma interpolação;
  \item avaliar em ciclo fechado, voltando a planear a cada segundo, com
    um critério de colisão baseado nas dimensões dos veículos;
  \item acrescentar variáveis do mapa e dos semáforos à previsão;
  \item aumentar o número de cenas e separar treino e teste por local,
    porque o mesmo local aparece em várias cenas;
  \item comparar o A* com outras estratégias de pesquisa e medir o número
    de estados expandidos.
\end{itemize}
